% Einheit 3 von 4 – Parameter und Transformationen (bank/trigonometrische-funktionen/e3.jsonl, zusammenbau v0.1)
%% TODO Zweigzeile Teil 1: was hier gelernt wird (Satz in Schülersprache aus dem Einheitstitel) – Einheit 3
\einheitenkopf[e3]{Einheit 3 von 4 · Parameter und Transformationen}
\zweigzeile{ab Kl. 10 · keine P10-Aufgabe}
\setcounter{aufgabe}{17}

%% TODO Titel: Ich-kann-Satz fehlt in der Bank (Platzhalter: Kettenname) – Nr. 18
%% TODO Anweisung: ein Satz über der ersten Teilaufgabe fehlt in der Bank – Nr. 18
\begin{aufgabe}{Parameter und Transformationen}
%% TODO \rechenplatz{n}: Schrittzahl des Grundfalls fehlt in der Bank (nur bei fester Schreibform, 2.3 b)
\begin{teile}
\teil Höher oder schmaler? Was hat sich bei g gegenüber f geändert?\\ \kreuz{die Höhe (a)}\kreuz{die Breite (b)}

\begin{ksys}[xmin=0,xmax=360,ymin=-2,ymax=2,xstep=45,ystep=0.5,ablesen]\funktion{sin(\x)}{f}\funktion{2*sin(\x)}{g}\end{ksys}
\teil y = 4 · sin(x) – wie groß ist die Amplitude? \leerfeld
\teil y = 5 · sin(x) – zwischen welchen Werten schwingt die Welle? \leerfeld
\teil Lies die Amplitude am Graphen ab.

\begin{ksys}[xmin=0,xmax=360,ymin=-3,ymax=3,xstep=45,ystep=0.5,ablesen]\funktion{2.5*sin(\x)}{f}\end{ksys}
$a = $ \leerfeld
\teil y = −2 · sin(x) – was ist anders als bei y = 2 · sin(x)? \leerfeld
\teil y = sin(3 · x) – wie lang ist die Periode? \leerfeld[°]
\teil Lies die Periode ab und berechne b in y = sin(b · x).

\begin{ksys}[xmin=0,xmax=360,ymin=-2,ymax=2,xstep=45,ystep=0.5,ablesen]\funktion{sin(4*\x)}{f}\end{ksys}
$b = $ \leerfeld
\teil Skizziere den Graphen von y = 2 · sin(3 · x) im ersten Vollkreis ab 0°.

\begin{ksys}[xmin=0,xmax=360,ymin=-3,ymax=3,xstep=30,ystep=1]\end{ksys}
\teil Stelle die Gleichung y = a · sin(b · x) zum Graphen auf: erst Amplitude, dann Periode, dann b.

\begin{ksys}[xmin=0,xmax=480,ymin=-5,ymax=5,xstep=60,ystep=1,ablesen]\funktion{4*sin(1.5*\x)}{f}\end{ksys}
$y = $ \leerfeld
\teil Welche Gleichung gehört zu g?\\ \kreuz{y = 2 · sin(x)}\\ \kreuz{y = sin(4 · x)}\\ \kreuz{y = 4 · sin(x)}

\begin{ksys}[xmin=0,xmax=360,ymin=-3,ymax=3,xstep=45,ystep=1,ablesen]\funktion{2*sin(\x)}{f}\funktion{sin(4*\x)}{g}\end{ksys}
\teil Der Faktor vor sin steigt von 1 auf 4. Beschreibe, was mit der Welle passiert.
\teil y = sin(x) + 2 – wohin ist die Welle verschoben, zwischen welchen Werten schwingt sie? \leerfeld
\teil Um wie viel Grad muss man y = sin(x) nach links schieben, damit y = cos(x) entsteht? Schreibe die Gleichung mit c. \leerfeld
\teil y = 2 · sin(4 · x) – erste Nullstelle nach 0° und erste Hochstelle? \leerfeld
\teil a) Stelle zum Graphen f die Gleichung y = a · sin(b · x) auf. b) Skizziere daneben den Graphen von y = 2 · sin(1,5 · x).

\begin{ksys}[xmin=0,xmax=360,ymin=-4,ymax=4,xstep=45,ystep=1,ablesen]\funktion{3.5*sin(3*\x)}{f}\end{ksys} \begin{ksys}[xmin=0,xmax=480,ymin=-3,ymax=3,xstep=60,ystep=1]\end{ksys}
\end{teile}
\end{aufgabe}

%% TODO Titel: Ich-kann-Satz fehlt in der Bank (Platzhalter: Kettenname) – Nr. 19
%% TODO Anweisung: ein Satz über der ersten Teilaufgabe fehlt in der Bank – Nr. 19
\begin{aufgabe}{zwei Graphen mit verschiedenen Parametern vergleichen}
\begin{teile}
\teil Vergleiche f und g im Bild: Was ist gleich, was ist verschieden?

\begin{ksys}[xmin=0,xmax=360,ymin=-3,ymax=3,xstep=45,ystep=1,ablesen]\funktion{2*sin(\x)}{f}\funktion{2*sin(3*\x)}{g}\end{ksys}
\end{teile}
\end{aufgabe}

%% TODO Titel: Ich-kann-Satz fehlt in der Bank (Platzhalter: Kettenname) – Nr. 20
%% TODO Anweisung: ein Satz über der ersten Teilaufgabe fehlt in der Bank – Nr. 20
\begin{aufgabe}{Parameter und Transformationen – Fehler finden, Begründen}
\begin{teile}
\teil Lisa schreibt zu y = sin(4 · x): Die Periode ist 1440°, denn b = 4 macht die Welle breiter. Finde den Fehler.
\teil Begründe, warum ein größeres b die Welle schmaler macht.
\end{teile}
\end{aufgabe}
